\chapter{第十一章：多智能体系统}

单个 Agent 很强，但它不一定适合所有问题。现实世界里，复杂任务往往天然具有分工：有人搜集信息，有人做分析，有人负责审稿，有人做最终决策。多智能体系统（multi-agent system）正是把这种组织结构搬进 AI 系统里。


从转行与面试角度看，多智能体不是``更酷的 demo''，而是你是否具备\textbf{系统拆分、协作协议设计、成本与延迟权衡、故障控制}能力的试金石。


\section{11.1 为什么需要多智能体}

单 Agent 常见瓶颈：


\begin{enumerate}[leftmargin=2em]
\item 工具集过大，选择困难；
\item 上下文过长，推理噪声增加；
\item 既要检索、又要写作、又要审校，角色冲突；
\item 单次失败影响整个任务；
\item 很难并行。
\end{enumerate}

而多 Agent 的核心收益有四个：


\begin{itemize}[leftmargin=2em]
\item \textbf{专业化（specialization）}：每个 Agent 只做一类事；
\item \textbf{并行化（parallelism）}：多个子任务同时执行；
\item \textbf{可替换（replaceability）}：某个 Agent 可独立迭代；
\item \textbf{可治理（governability）}：更容易限制权限和审计。
\end{itemize}

但要牢记一句话：\textbf{多 Agent 不是免费午餐}。你把一个模型调用拆成 3 个 Agent，成本、延迟、协调复杂度通常会同步上升。


\section{11.2 五种多智能体通信模式}

\subsection{11.2.1 Sequential Chain（顺序链）}

最简单的模式：Agent A 输出交给 Agent B，再交给 Agent C。


\begin{codeblock}
User Task
   |
   v
[Researcher] --> [Writer] --> [Reviewer] --> Final Output
\end{codeblock}

适用场景：


\begin{itemize}[leftmargin=2em]
\item 数据清洗 -> 生成报告 -> 质量检查；
\item 检索 -> 总结 -> 翻译；
\item 需求理解 -> 代码生成 -> 测试建议。
\end{itemize}

优点：


\begin{itemize}[leftmargin=2em]
\item 结构清晰；
\item 调试简单；
\item 每步职责明确。
\end{itemize}

缺点：


\begin{itemize}[leftmargin=2em]
\item 串行延迟高；
\item 上游错误会层层传递；
\item 中间格式必须设计好。
\end{itemize}

\subsection{11.2.2 Parallel Fan-out（并行扇出）}

由一个 orchestrator（编排器）把任务分发给多个 Agent 并行执行，再聚合结果。


\begin{codeblock}
                 +--> [Agent A: 法规检索] ---+
User Task --> Orchestrator --> [Agent B: 市场数据] --+--> Synthesizer
                 +--> [Agent C: 技术趋势] ---+
\end{codeblock}

适合：


\begin{itemize}[leftmargin=2em]
\item 多来源调研；
\item 多城市、多国家对比；
\item 多候选方案评估。
\end{itemize}

它通常能把 15 秒串行任务压缩到 5\textasciitilde{}6 秒，但代价是同时消耗更多 token 与 API quota。


\subsection{11.2.3 Hierarchical（层级式）}

Supervisor Agent（主管 Agent）负责规划和分配，Specialist Agents（专家 Agent）负责执行。


\begin{codeblock}
                 +--> [Search Specialist]
[Supervisor] ----+--> [Coding Specialist]
                 +--> [Review Specialist]
\end{codeblock}

这是企业里最常见、也最像组织结构的模式。主管 Agent 不一定自己做事，而是：


\begin{itemize}[leftmargin=2em]
\item 拆任务；
\item 指派角色；
\item 跟踪进度；
\item 决定是否重试或改派。
\end{itemize}

\subsection{11.2.4 Debate / Adversarial（辩论 / 对抗）}

多个 Agent 持不同立场，最后由裁判或聚合器综合。


例如：


\begin{itemize}[leftmargin=2em]
\item Agent A：支持采用 LangGraph；
\item Agent B：反对，强调维护成本；
\item Judge：根据证据做综合结论。
\end{itemize}

这种模式适合：


\begin{itemize}[leftmargin=2em]
\item 决策质量要求高；
\item 需要显式暴露争议点；
\item 防止单模型过早收敛。
\end{itemize}

但成本高，且若没有证据约束，很容易变成``高 token 的空谈''。


\subsection{11.2.5 Swarm（群体式）}

Swarm 指没有强中心控制，Agent 根据能力和局部状态自组织完成任务。它更接近分布式系统而不是简单工作流。


适合：


\begin{itemize}[leftmargin=2em]
\item 大规模模拟；
\item 动态任务市场；
\item 去中心化协作实验。
\end{itemize}

工程落地难点最大，包括：


\begin{itemize}[leftmargin=2em]
\item 发现机制；
\item 冲突解决；
\item 重复劳动抑制；
\item 全局终止条件。
\end{itemize}

对转行求职者来说，理解概念即可，真正项目中最常见还是顺序链、并行扇出和层级式。


\section{11.3 编排设计（Orchestration Design）}

\subsection{11.3.1 中央编排器模式}

中央编排器模式中，所有任务调度由一个核心组件控制：


\begin{codeblock}
                   +-------------------+
                   |   Orchestrator    |
                   +---------+---------+
                             |
        +--------------------+--------------------+
        |                    |                    |
        v                    v                    v
   [Agent A]            [Agent B]            [Agent C]
\end{codeblock}

优点：


\begin{itemize}[leftmargin=2em]
\item 逻辑集中，便于监控；
\item 权限控制容易；
\item 失败恢复可统一处理。
\end{itemize}

缺点：


\begin{itemize}[leftmargin=2em]
\item 单点瓶颈；
\item 中心调度可能成为延迟热点；
\item 可扩展性受限。
\end{itemize}

\subsection{11.3.2 去中心化 Peer-to-Peer}

Agent 之间直接通信，没有唯一中央调度者。


优点：


\begin{itemize}[leftmargin=2em]
\item 灵活；
\item 可扩展；
\item 某种程度上更鲁棒。
\end{itemize}

缺点：


\begin{itemize}[leftmargin=2em]
\item 协议设计复杂；
\item 状态一致性难；
\item 调试非常痛苦。
\end{itemize}

在面试中，如果对方问``生产里你更偏好哪种？''大多数场景下回答\textbf{中心编排优先}是更稳妥的。


\subsection{11.3.3 Blackboard / Shared Memory（黑板模式）}

黑板模式借鉴经典 AI：所有 Agent 从共享内存空间读取任务与事实，再写回结果。


\begin{codeblock}
           +------------------------+
           |   Shared Blackboard    |
           | task / evidence / logs |
           +-----------+------------+
                       ^
          +------------+------------+
          |            |            |
      [Agent A]    [Agent B]    [Agent C]
\end{codeblock}

适合：


\begin{itemize}[leftmargin=2em]
\item 需要共享证据；
\item 多 Agent 逐步完善同一结果；
\item 长任务可断点续跑。
\end{itemize}

共享黑板常落地为：


\begin{itemize}[leftmargin=2em]
\item Redis；
\item Postgres；
\item 文档数据库；
\item 事件总线 + 状态存储。
\end{itemize}

\section{11.4 多智能体框架对比}

\begingroup
\scriptsize
\setlength{\tabcolsep}{3pt}
\begin{longtable}{|p{0.19\textwidth}|p{0.19\textwidth}|p{0.19\textwidth}|p{0.19\textwidth}|p{0.19\textwidth}|}
\hline
Framework & Pattern & Strengths & Weaknesses & Best For \\ \hline
CrewAI & Role-based sequential / hierarchical & 概念直观、上手快、适合 demo & 状态控制不够细、复杂图编排一般 & 内容生产、轻量团队协作 \\ \hline
AutoGen & Conversational multi-agent & 对话式协作灵活、适合研究型实验 & 对生产治理要求高、token 成本易失控 & 原型验证、agent chat \\ \hline
LangGraph & Graph/state machine & 状态显式、可控、适合生产级复杂流程 & 学习曲线更陡 & 可恢复工作流、复杂编排 \\ \hline
Claude Code Coordinator & Supervisor + specialists & 工程任务分工自然、适合代码/CLI 任务 & 依赖具体执行环境 & 开发代理、命令执行协作 \\ \hline
OpenAI Swarm & Lightweight handoff & handoff 简单、概念清晰 & 更偏实验与轻编排 & 小型多 Agent 切换 \\ \hline
\end{longtable}
\endgroup

面试时最好不要只背框架名字，而要说明你选型的理由。例如：


\begin{quote}
如果任务高度状态化、需要 checkpoint 和人工审核，我更倾向 LangGraph；如果只是角色分工清晰的内容任务，CrewAI 足够。\\
\end{quote}

\section{11.5 实际难题：多 Agent 比单 Agent 更容易出事}

\subsection{11.5.1 成本爆炸}

假设一个单 Agent 任务平均 8k input token + 2k output token。如果拆成 3 个 Agent：


\begin{itemize}[leftmargin=2em]
\item Researcher：6k + 1k
\item Writer：8k + 2k
\item Reviewer：7k + 1k
\item Orchestrator：3k + 0.5k
\end{itemize}

总量可能从 10k 变成 28.5k token，直接接近 3 倍。你必须问自己：质量提升是否值得这 3 倍成本？


\subsection{11.5.2 协调失败与死锁}

常见失败：


\begin{itemize}[leftmargin=2em]
\item Agent A 等 Agent B 结果，B 又在等 A；
\item 多个 Agent 重复做同一件事；
\item 上下游输出格式不兼容；
\item 一个 Agent 无限请求``再给我更多上下文''。
\end{itemize}

解决思路：


\begin{itemize}[leftmargin=2em]
\item 显式任务状态机；
\item 每个 Agent 最大轮数限制；
\item 明确定义输入输出 schema；
\item central timeout + cancellation。
\end{itemize}

\subsection{11.5.3 上下文共享}

共享太少：


\begin{itemize}[leftmargin=2em]
\item Agent 缺背景，做错事。
\end{itemize}

共享太多：


\begin{itemize}[leftmargin=2em]
\item token 成本高；
\item 隐私边界模糊；
\item 噪声干扰判断。
\end{itemize}

实践建议：


\begin{itemize}[leftmargin=2em]
\item 共享摘要，不共享全量；
\item 共享证据链接，不共享全文；
\item 共享结构化 facts，不共享冗长聊天。
\end{itemize}

\subsection{11.5.4 调试困难}

单 Agent 调试只看一条 trace，多 Agent 要看：


\begin{itemize}[leftmargin=2em]
\item 谁发给谁；
\item 每个 Agent 看到了什么上下文；
\item 是谁做的最终决策；
\item 某条错误信息是否被正确传播。
\end{itemize}

所以必须记录：


\begin{itemize}[leftmargin=2em]
\item agent\_id
\item parent\_task\_id
\item message\_id
\item input snapshot
\item output snapshot
\item latency
\item token usage
\end{itemize}

\subsection{11.5.5 延迟管理}

多 Agent 最大的用户体验问题常常不是错，而是慢。


优化手段包括：


\begin{itemize}[leftmargin=2em]
\item 并行优先；
\item 设定最快足够策略（good-enough threshold）；
\item 提前流式返回中间进度；
\item 对稳定中间结果做缓存；
\item 对 Reviewer 只检查高风险段落，而非全文重审。
\end{itemize}

\section{11.6 示例：构建一个 3 Agent 研究团队}

目标：给用户生成一份``某技术主题研究报告''。


角色分工：


\begin{enumerate}[leftmargin=2em]
\item \textbf{Researcher Agent}：搜索、收集事实；
\item \textbf{Writer Agent}：把资料整合成报告；
\item \textbf{Reviewer Agent}：检查准确性、缺口和表达质量。
\end{enumerate}

\subsection{11.6.1 通信架构图}

\begin{codeblock}
                            +----------------------+
                            |   User / Client      |
                            +----------+-----------+
                                       |
                                       v
                            +----------+-----------+
                            |   Orchestrator       |
                            | task_id / state /    |
                            | retries / timeouts   |
                            +----+-----------+-----+
                                 |           |
                                 |           |
                                 v           |
                      +----------+-----+     |
                      | Researcher     |-----+
                      | search, fetch  |     |
                      +----------+-----+     |
                                 | evidence  |
                                 v           |
                      +----------+-----+     |
                      | Writer         |-----+
                      | draft report   |     |
                      +----------+-----+     |
                                 | draft     |
                                 v           |
                      +----------+-----+     |
                      | Reviewer       |-----+
                      | critique, fix  |
                      +----------------+
\end{codeblock}

\subsection{11.6.2 完整 Python 示例}

下面给一个最小可运行版本。为避免依赖真实模型 API，这里用简单函数模拟三个 Agent 的行为，重点展示 orchestration 结构。


\begin{codeblock}
from __future__ import annotations

from dataclasses import dataclass, field
from typing import Dict, List


def fake_search(topic: str) -> List[str]:
    corpus = {
        "mcp": [
            "MCP uses JSON-RPC 2.0 as the base protocol.",
            "MCP supports tools, resources, prompts and sampling.",
            "MCP commonly runs over stdio, HTTP+SSE or WebSocket."
        ],
        "a2a": [
            "A2A focuses on agent-to-agent task coordination.",
            "Agent Cards advertise capabilities and auth requirements.",
            "OAuth 2.1 is commonly used for authentication."
        ],
    }
    return corpus.get(topic.lower(), [f"No data found for {topic}"])


@dataclass
class AgentMessage:
    sender: str
    content: str


class ResearcherAgent:
    name = "researcher"

    def run(self, topic: str) -> Dict[str, List[str]]:
        facts = fake_search(topic)
        return {"topic": topic, "facts": facts}


class WriterAgent:
    name = "writer"

    def run(self, research_result: Dict[str, List[str]]) -> str:
        bullet_points = "\n".join(f"- {x}" for x in research_result["facts"])
        return f"# Research Report: {research_result['topic']}\n\n## Key Findings\n{bullet_points}\n"


class ReviewerAgent:
    name = "reviewer"

    def run(self, draft: str) -> Dict[str, str]:
        issues = []
        if "No data found" in draft:
            issues.append("缺少有效证据")
        if "Key Findings" not in draft:
            issues.append("结构不完整")
        if issues:
            return {"approved": "false", "feedback": "；".join(issues)}
        return {"approved": "true", "feedback": "报告结构完整，已包含主要事实。"}


@dataclass
class Orchestrator:
    researcher: ResearcherAgent = field(default_factory=ResearcherAgent)
    writer: WriterAgent = field(default_factory=WriterAgent)
    reviewer: ReviewerAgent = field(default_factory=ReviewerAgent)
    logs: List[AgentMessage] = field(default_factory=list)

    def execute(self, topic: str) -> Dict[str, str]:
        research = self.researcher.run(topic)
        self.logs.append(AgentMessage("researcher", str(research)))

        draft = self.writer.run(research)
        self.logs.append(AgentMessage("writer", draft))

        review = self.reviewer.run(draft)
        self.logs.append(AgentMessage("reviewer", str(review)))

        if review["approved"] == "true":
            return {"status": "done", "report": draft, "review": review["feedback"]}

        improved_draft = draft + f"\n## Reviewer Feedback\n{review['feedback']}\n"
        return {"status": "needs_revision", "report": improved_draft, "review": review["feedback"]}


if __name__ == "__main__":
    orchestrator = Orchestrator()
    result = orchestrator.execute("mcp")
    print(result["status"])
    print(result["report"])
    print(result["review"])

    print("\n=== Logs ===")
    for log in orchestrator.logs:
        print(f"[{log.sender}] {log.content}")
\end{codeblock}

\subsection{11.6.3 如何升级到生产版本}

生产化时，建议加入：


\begin{itemize}[leftmargin=2em]
\item 每个 Agent 单独 system prompt；
\item 工具权限差异化；
\item 中间结果 schema；
\item 任务状态表；
\item 审稿打分标准；
\item 人工接管入口；
\item agent-level timeout；
\item streaming progress。
\end{itemize}

\section{11.7 什么时候不应该用多 Agent}

以下场景用单 Agent 更合适：


\begin{itemize}[leftmargin=2em]
\item 任务只需 1\textasciitilde{}3 次工具调用；
\item 所有步骤强顺序且不复杂；
\item 上下文共享极多，拆分后反而重复传输；
\item 延迟预算极小，比如必须 1 秒内返回；
\item 成本极度敏感。
\end{itemize}

一个很成熟的工程判断是：\textbf{先用单 Agent 跑通，再在瓶颈处引入多 Agent，而不是从第一天就把系统拆成``AI 微服务集群''。}


\section{11.8 任务应该如何拆给多个 Agent}

多 Agent 最大的设计问题，不是`\texttt{能不能再加一个 Agent''，而是}`任务边界怎么划分才合理''。一个实用判断标准是看子任务是否同时满足四个条件：


\begin{enumerate}[leftmargin=2em]
\item \textbf{输入明确}：拿到什么信息才能开工；
\item \textbf{输出明确}：产出是什么格式；
\item \textbf{能力独立}：是否需要独立工具集或独立角色提示；
\item \textbf{可验证}：结果能否被下游或审稿者检查。
\end{enumerate}

例如`\texttt{研究 + 写作 + 审核''适合拆，因为三者目标不同、输出也不同；但}`先写第一段，再写第二段，再写第三段''往往不值得拆，因为强依赖同一上下文，拆开只会增加拼接成本。


可以用下面这张表快速判断：


\begingroup
\scriptsize
\setlength{\tabcolsep}{3pt}
\begin{longtable}{|p{0.31\textwidth}|p{0.31\textwidth}|p{0.31\textwidth}|}
\hline
问题 & 若回答是``是'' & 更倾向多 Agent \\ \hline
子任务是否可并行 & 多地调研、多个数据源 & 是 \\ \hline
是否需要不同角色视角 & 研究、写作、审核 & 是 \\ \hline
是否需要不同权限 & 只读检索 vs 写入系统 & 是 \\ \hline
子任务之间是否高度耦合 & 共享全部上下文 & 否 \\ \hline
\end{longtable}
\endgroup

工程上一个很重要的原则是：\textbf{按能力拆，不按想象中的``团队角色''硬拆}。否则很容易出现 5 个 Agent 名字很酷，但只有 1 个在真正做事。


\section{11.9 多 Agent 之间共享什么上下文}

多 Agent 协作最忌讳``全量转发全部历史''。更推荐共享结构化上下文包：


\begin{codeblock}
{
  "task_id": "task-2026-0612-001",
  "goal": "生成 MCP 与 A2A 对比报告",
  "constraints": ["必须给出协议层差异", "需要带代码示例"],
  "evidence": [
    {"source": "doc-1", "fact": "MCP 基于 JSON-RPC 2.0"},
    {"source": "doc-2", "fact": "A2A 强调 Agent Card 与任务生命周期"}
  ],
  "budget": {
    "max_tokens": 6000,
    "deadline_seconds": 30
  }
}
\end{codeblock}

这样的好处有三点：


\begin{enumerate}[leftmargin=2em]
\item 下游 Agent 看到的是``任务对象''，不是一团聊天记录；
\item 编排器可插入成本、时间、权限约束；
\item 审稿 Agent 可以只看证据和草稿，不必看全部搜索过程。
\end{enumerate}

在生产系统中，通常会把共享内容分层：


\begin{itemize}[leftmargin=2em]
\item \textbf{必须共享}：目标、约束、证据、草稿版本；
\item \textbf{可选共享}：中间推理摘要、失败原因；
\item \textbf{不共享}：冗长内部思考、无关聊天、敏感凭证。
\end{itemize}

\section{11.10 生产治理：避免``Agent 会议开不完''}

多 Agent 常见坏味道是：每个 Agent 都很能说，但系统不收敛。解决方法不是写更长 prompt，而是做治理。


\subsection{一、给每个 Agent 设退出条件}

例如：


\begin{itemize}[leftmargin=2em]
\item Researcher 最多检索 5 个来源；
\item Writer 最多修订 2 轮；
\item Reviewer 若无致命问题则必须通过，不允许无限挑小毛病。
\end{itemize}

\subsection{二、给 Orchestrator 设总预算}

总预算包括：


\begin{itemize}[leftmargin=2em]
\item 总 token；
\item 总 wall-clock 时间；
\item 总调用轮数；
\item 总并行度。
\end{itemize}

如果预算快用完，编排器应主动降级，例如减少 reviewer 深度、只保留前 3 个证据源。


\subsection{三、日志必须可回放}

你需要能回答：


\begin{itemize}[leftmargin=2em]
\item 哪个 Agent 首先引入了错误事实；
\item 哪个 Agent 忽略了关键约束；
\item 某个任务为什么卡了 40 秒；
\item 为什么这次成本是上次的 2 倍。
\end{itemize}

因此每条通信都应带：


\begin{itemize}[leftmargin=2em]
\item \texttt{task\_id}
\item \texttt{agent\_id}
\item \texttt{parent\_agent\_id}
\item \texttt{message\_type}
\item \texttt{created\_at}
\item \texttt{latency\_ms}
\item \texttt{token\_usage}
\end{itemize}

\subsection{四、对人类可解释}

多 Agent 系统上线后，业务方最常问的不是`\texttt{用了几个模型''，而是}`为什么这个结论是这么来的''。因此最好保留：


\begin{itemize}[leftmargin=2em]
\item 证据链；
\item 角色分工；
\item 审核意见；
\item 最终裁决依据。
\end{itemize}

\section{11.11 什么时候需要 Debate，什么时候不需要}

辩论式模式很吸引人，因为它看起来像``让两个专家互相挑错''。但不是所有任务都值得这样做。适合 Debate 的场景通常有两个特征：


\begin{enumerate}[leftmargin=2em]
\item 任务存在多种合理方案；
\item 你能给出明确裁判标准。
\end{enumerate}

例如`\texttt{选 LangGraph 还是 AutoGen 做下一代平台''的确值得让两个 Agent 站在不同立场辩论；但}`把数据库连接串从配置里读出来''这类任务就完全没必要辩论。


如果没有裁判标准，Debate 很容易退化成：


\begin{itemize}[leftmargin=2em]
\item 双方重复相同事实；
\item 消耗双倍 token；
\item 最后仍需要人工判断。
\end{itemize}

所以实战里，Debate 应该是高价值任务的增强模块，而不是默认配置。


\section{11.12 一个多 Agent 平台的运行指标}

如果单 Agent 关注的是工具成功率和 loop 次数，多 Agent 还要额外看协作指标：


\begingroup
\scriptsize
\setlength{\tabcolsep}{3pt}
\begin{longtable}{|p{0.47\textwidth}|p{0.47\textwidth}|}
\hline
指标 & 含义 \\ \hline
delegation\_rate & 任务被拆分与委派的比例 \\ \hline
parallelism\_factor & 平均并发子任务数 \\ \hline
handoff\_success\_rate & Agent 之间交接成功率 \\ \hline
revision\_rounds & 草稿被 reviewer 打回的次数 \\ \hline
duplicate\_work\_rate & 多个 Agent 重复劳动比例 \\ \hline
cost\_per\_completed\_task & 完成一个任务的总成本 \\ \hline
\end{longtable}
\endgroup

其中 \texttt{duplicate\_work\_rate} 非常关键。很多多 Agent 系统表面看起来``很勤奋''，实际上 3 个 Agent 在查同样的资料，只是用不同表述重复劳动。解决这个问题，往往要靠共享证据池和任务去重键，而不是靠更强模型。


\section{11.13 从单 Agent 迁移到多 Agent 的推荐路径}

最安全的方式不是``一步拆成 5 个 Agent''，而是逐步迁移：


\subsection{阶段一：单 Agent + 多工具}

先把任务跑通，搞清楚真正瓶颈是检索、写作、审校还是执行。


\subsection{阶段二：拆出第一个专家 Agent}

通常先拆检索或审稿，因为它们职责边界最清晰。比如保留主 Agent 负责对话，把``资料收集''交给 Research Agent。


\subsection{阶段三：引入编排器}

只有当你确定存在多角色、多阶段、多预算控制需求时，再引入 Supervisor 或 Orchestrator。否则编排器本身会成为额外复杂度来源。


\subsection{阶段四：协议化与治理}

当 Agent 数量增多后，再做消息 schema、共享状态、权限和可观测性统一。


这种迁移思路在面试里很有说服力，因为它体现了你不会为了架构而架构，而是先验证价值，再增加复杂度。


\section{11.14 一个常见失败案例}

某团队想做``自动写周报系统''，一开始就设计了选题 Agent、搜集 Agent、写作 Agent、润色 Agent、审稿 Agent、排版 Agent 六个角色。结果上线后发现：


\begin{itemize}[leftmargin=2em]
\item 任务本身并不复杂；
\item 六个 Agent 共享几乎同一份上下文；
\item 每轮都要传递大段草稿；
\item 延迟从 6 秒变成 28 秒；
\item 成本接近单 Agent 的 4 倍。
\end{itemize}

最后他们把系统收缩为：Researcher + Writer + Reviewer 三个角色，成本与时延显著下降，质量变化却不大。


这个案例说明，多 Agent 的核心不是``角色越多越专业''，而是\textbf{边界越清晰越高效}。


\section{11.15 一个适合面试展示的系统设计回答}

如果面试官让你设计``一个自动生成行业研究报告的多 Agent 系统''，你可以按下面顺序回答：


第一，定义角色：


\begin{itemize}[leftmargin=2em]
\item Researcher：负责搜索、抓取、提取事实；
\item Analyst：负责把事实转成结构化比较；
\item Writer：负责成稿；
\item Reviewer：负责事实核验和表达质量；
\item Orchestrator：负责分配任务、汇总状态、控制预算。
\end{itemize}

第二，定义通信协议：


\begin{itemize}[leftmargin=2em]
\item Researcher 向共享黑板写证据条目；
\item Analyst 读取证据并生成对比表；
\item Writer 只读对比表和证据摘要；
\item Reviewer 输出问题清单与通过/不通过结果。
\end{itemize}

第三，定义状态：


\begin{itemize}[leftmargin=2em]
\item 任务级状态：创建、执行中、待审稿、完成；
\item Agent 级状态：空闲、运行中、阻塞、失败；
\item 资源级状态：预算剩余、超时、证据数量。
\end{itemize}

第四，定义治理：


\begin{itemize}[leftmargin=2em]
\item 每个 Agent 最大轮次；
\item 最大并发数；
\item 证据最少来源数量，比如至少 3 个独立来源；
\item 总成本超过阈值时降级为``短报告模式''。
\end{itemize}

第五，定义可观测性：


\begin{itemize}[leftmargin=2em]
\item 每个 handoff 的延迟；
\item 每个 Agent 的 token 消耗；
\item 草稿被 reviewer 打回的次数；
\item 最终报告的引用覆盖率。
\end{itemize}

这种答题方式的优势是，它让``多 Agent''听起来像真正的分布式协作系统，而不只是几个 Prompt 拼起来。


\section{11.16 为什么多 Agent 特别适合传统软件工程师转型}

因为它几乎把所有你熟悉的工程能力都重新用了一遍：


\begin{itemize}[leftmargin=2em]
\item 任务拆分类似服务边界设计；
\item Agent 间通信类似 RPC 或消息队列；
\item 共享黑板类似数据库或事件总线；
\item 重试、幂等、超时、取消都是分布式系统老问题；
\item 成本、延迟、吞吐量权衡也和后端系统非常像。
\end{itemize}

所以你不必把多 Agent 看成一个神秘的新世界。它本质上是：在经典系统设计问题上，再加一层由 LLM 驱动的决策与语言接口。这也是为什么面试官特别喜欢在这一章考察候选人的工程感觉。


\section{11.17 一个落地判断标准}

如果你在项目里犹豫``要不要上多 Agent''，可以问自己三个问题：

第一，这个任务是否真的存在可独立优化的子角色？

第二，拆开后是否能换来并行、质量或权限上的明显收益？

第三，新增的协调成本是否小于收益？


只要这三个问题里有两个回答是否定的，就先别急着拆。这个判断标准非常朴素，但在真实项目里极其有用。


也正因为如此，优秀的多 Agent 设计往往显得`\texttt{克制''：角色不多、协议清楚、边界稳定、日志完整。它追求的是系统收益最大化，而不是视觉上看起来像一个}`很热闹的 AI 团队''。

真正优秀的方案，常常是把复杂度精确地放在最有价值的环节，而不是平均分散到每个 Agent 身上。

如果你能在设计中始终回答``为什么必须拆、拆开后收益是什么、谁来协调失败''，那你的多 Agent 架构通常就不会走偏。

否则，它很容易从协作系统退化成昂贵而混乱的消息风暴。

从系统设计视角看，多 Agent 的难点和微服务并没有本质区别：服务拆分之后，真正的复杂度会转移到通信、状态一致性、失败恢复和观测层。差别只在于，这里的``服务''既有确定性程序，也有带概率性的模型决策模块。所以，越是复杂的多 Agent 系统，越要在接口、状态机和回放能力上做硬约束，而不是把希望寄托在 Prompt 足够长、模型足够聪明上。

这也是为什么越成熟的团队，越会把多 Agent 当成系统工程题，而不是单纯的 Prompt 题。

一旦把这个认知立住，你的架构判断会稳很多。

你会更愿意先算清收益，再决定要不要拆分角色、增加链路和引入新的协作状态。

这是一种非常工程化的思维：先确认瓶颈，再增加结构；先证明价值，再引入复杂度。对转型中的软件工程师来说，这正是你相对纯算法背景候选人的优势所在。

多 Agent 不是`\texttt{会不会搭''的问题，更是}`值不值得搭、搭完如何收敛''的问题。

只有把收敛性、预算和可观测性一起设计进去，多 Agent 才会从概念图变成真正能上线的系统。

换句话说，多 Agent 的上限不由角色数量决定，而由协作质量决定。


\section{本章要点}

\begin{itemize}[leftmargin=2em]
\item 多智能体系统适合复杂任务分工、并行处理和角色专业化。
\item 常见模式包括顺序链、并行扇出、层级式、辩论式和群体式。
\item 生产中最常见且最稳妥的是中央编排 + 专家 Agent 的层级式设计。
\item 多 Agent 的主要代价是成本上升、协调失败、上下文共享困难、调试复杂和延迟增加。
\item 框架选型要看状态管理需求、可恢复性、成本控制和开发复杂度，而不是只看流行度。
\item 一个好的多 Agent 设计，必须先定义协议、状态、边界，再讨论模型与 prompt。
\end{itemize}

\section{延伸阅读}

\begin{enumerate}[leftmargin=2em]
\item AutoGen 官方论文与文档
\item CrewAI 官方文档
\item LangGraph 多 Agent / state graph 设计指南
\item OpenAI Swarm 示例仓库
\item 关于 blackboard architecture 与 distributed systems coordination 的经典资料
\end{enumerate}
